\chapter{PROSES PEMBUATAN PRODUK}

\section{Lingkungan Implementasi}

Implementasi dilakukan pada klaster Kubernetes yang dikelola oleh Rancher dengan spesifikasi sebagai berikut:

\begin{itemize}
    \item \textbf{Sistem Operasi:} Linux (Ubuntu 22.04 LTS)
    \item \textbf{Container Runtime:} containerd
    \item \textbf{Kubernetes:} v1.28
    \item \textbf{ArgoCD:} v2.9
    \item \textbf{Kustomize:} v5.0
\end{itemize}

\section{Implementasi GitOps Workflow}

\subsection{Struktur Repositori}

Repositori Git disusun dengan struktur yang terorganisir:

\begin{lstlisting}[language=bash, caption={Struktur Repositori GitOps}]
argocd/
  apps/                 # ArgoCD Application manifests
  bootstrap/            # Bootstrap configurations
k8s/
  base/                 # Base Kubernetes manifests
  overlays/             # Environment-specific overlays
  secrets/              # Sealed Secrets (encrypted)
scripts/
  verify-infra.sh       # Infrastructure verification
  bootstrap-infra.sh    # Infrastructure bootstrap
\end{lstlisting}

\subsection{Konfigurasi ArgoCD Application}

Setiap komponen didefinisikan sebagai ArgoCD Application:

\begin{lstlisting}[language=yaml, caption={Contoh ArgoCD Application}]
apiVersion: argoproj.io/v1alpha1
kind: Application
metadata:
  name: invenio-bootstrap
  namespace: argocd
spec:
  project: default
  source:
    repoURL: https://github.com/vityasyyy/invenio-rdm-gitops
    targetRevision: HEAD
    path: k8s/overlays/production
  destination:
    server: https://kubernetes.default.svc
    namespace: invenio
  syncPolicy:
    automated:
      prune: true
      selfHeal: true
\end{lstlisting}

\section{Implementasi Keamanan}

\subsection{Network Policies}

Network Policies diimplementasikan dengan pendekatan default-deny:

\begin{lstlisting}[language=yaml, caption={Contoh Network Policy}]
apiVersion: networking.k8s.io/v1
kind: NetworkPolicy
metadata:
  name: invenio-network-policy
  namespace: invenio
spec:
  podSelector:
    matchLabels:
      app: invenio
  policyTypes:
  - Ingress
  - Egress
  ingress:
  - from:
    - namespaceSelector:
        matchLabels:
          name: traefik
    ports:
    - protocol: TCP
      port: 5000
\end{lstlisting}

\subsection{Sealed Secrets}

Manajemen rahasia diimplementasikan menggunakan Sealed Secrets.

\section{Implementasi Akses Eksternal}

Cloudflare Tunnel diimplementasikan sebagai DaemonSet untuk akses eksternal yang aman dengan pendekatan outbound-only.
